\documentclass[11pt]{article}

\usepackage[margin=1.1in]{geometry}
\usepackage{amsmath}
\usepackage{amssymb}
\usepackage{booktabs}
\usepackage{microtype}
\usepackage{xcolor}
\usepackage{graphicx}
\usepackage{caption}
\usepackage{array}
\usepackage[colorlinks=true,linkcolor=blue!50!black,citecolor=blue!50!black,urlcolor=blue!50!black]{hyperref}

\hypersetup{
  pdftitle={The Contained Book: An Absorbing Resource Measured to Be Worth Nothing},
  pdfauthor={Allen Hsieh}
}

\newtheorem{theorem}{Theorem}
\newtheorem{corollary}{Corollary}
\newtheorem{definition}{Definition}
\newenvironment{proofenv}{\par\noindent\textit{Proof.}\ }{\hfill$\square$\par\medskip}

\newcommand{\code}[1]{\texttt{#1}}

\title{The Contained Book:\\An Absorbing Resource Measured to Be Worth Nothing}
\author{Allen Hsieh\\[2pt]{\small FishAI project}}
\date{August 2026}

\begin{document}

\maketitle

\begin{abstract}
\noindent
In the six-player partnership game Literature (Canadian Fish), a half-suit whose
six cards all sit with one team is an \emph{absorbing state}: the opponents can
neither ask into it nor take it by declaring, so leaving it unclaimed is a
resource rather than an oversight. The resource is concrete --- a holder may
deliberately ask into the set, which is a guaranteed miss, and a miss hands the
turn to an opponent of the asker's choosing without moving a card. This yields a
renewable, aimable, nearly information-free turn-pass that only the containing
team can play and that claiming the set destroys. We prove the absorbing
property from the rules of both rule sets the FishAI engine supports ---
correcting, along the way, the project's own earlier claim that it failed under
the 48-card baseline --- derive a trigger inequality that prices the move in
cards against the ordinary ask it would displace, implement it as a policy
option with a strict no-convention discipline, and measure it at three scales:
per-decision mirrors, 1{,}200 duplicate-deal pairs per style, and a full
$36$-cell $\times$ 4{,}300-pair round-robin re-run. The mechanism fires, fires
for exactly the reasons the derivation gives, aims where the style asks, and
reuses one card so that 95--98\% of its uses publish nothing new. Its measured
value is zero: no style moves two standard errors from $0.5$ in the paired
mirror, and the full matrix re-run leaves the verdict, the ranking, and the
cyclic structure unchanged, with the heaviest intended user shifting by
$-0.0018 \pm 0.0014$. The arithmetic says why: at the roster's measured hit
rate a conceded turn is worth about one card, so the entire spread between the
best and worst opponent to concede to is about one card. We report this as a
model negative result --- about the move at this level of play, not about the
theorem.
\end{abstract}

\section{Introduction}
\label{sec:intro}

Negative results are cheap to assert and expensive to earn. This paper earns
one. It concerns a single strategic mechanism in Literature (Canadian Fish)
under the FishAI project's \code{us54} dialect: the \emph{contained book}, a
half-suit wholly held by one team. The mechanism has three parts, each of which
survives scrutiny --- and a value, which does not.

The parts. First, a theorem: containment is absorbing. No opponent can legally
ask into a contained book, and an opponent who declares it hands it to the
holders, so the position can only be ended by the holders themselves
(Section~\ref{sec:theorem}). Second, a move: a holder may ask into the book
anyway. The ask is legal, cannot hit, and on the miss the turn passes to
whichever opponent the asker named --- a renewable, aimable turn-pass that
costs at most one card of hand information over its lifetime if one card is
reused (Section~\ref{sec:move}). Third, a price: a derived inequality, with
every term measured off public data in units of cards, that says when the pass
beats the ordinary ask it would displace (Section~\ref{sec:trigger}).

The value. Implemented with that pricing and handed to a nine-style roster, the
move is worth nothing. Across 1{,}200 duplicate-deal pairs per style, no style
equipped with the move beats its own unequipped twin by two standard errors
(Section~\ref{sec:mirror}). Across a full-precision re-run of the project's
$36$-cell round-robin \cite{v05} --- same 4{,}300 seeds, same rotation, the policy change
and nothing else --- the verdict, the ranking, and the cyclic structure are all
unchanged (Section~\ref{sec:matrix}). The measured value of the contained-book
turn-pass, as a policy option under \code{us54} against the roster it was built
for, is zero \cite{styles}.

Two features make this negative result worth writing up rather than filing
away. The first is that every stage of it is \emph{measured rather than
asserted}, and twice the measurement corrected the project's own prose: the
handoff's claim that the move is information-free (it is not, but the
qualification makes the move better --- Section~\ref{sec:move}); and an earlier
draft's claim that containment is not absorbing under the 48-card baseline
rules (it is --- Section~\ref{sec:correction}). The second is that the
mechanism demonstrably works. It fires at the derived rate, refuses in the
derived cases, and never leaks --- so its worthlessness is a finding about the
game at this level of play, not a bug report. The arithmetic bounding the
effect at about one card (Section~\ref{sec:small}) tells us where value would
have to come from if it exists at all: the declare policy, which is untouched
here (Section~\ref{sec:open}).

\section{The rules the result rests on}
\label{sec:rules}

Literature is played by six players in two teams of three (seats $0,2,4$
versus $1,3,5$) \cite{pagat}. Under \code{us54} \cite{rules54} the deck is 54 cards ---
the standard 52 plus two distinguishable jokers --- partitioned into nine
\emph{books} (half-suits) of six cards: per suit, LOW $=$ 2--7 and HIGH $=$
9--A, plus an EIGHTS book of the four 8s and both jokers. Nine cards are dealt
to each player; the game ends the moment either team has been awarded five
books. The project's shipped baseline, \code{pagat48} \cite{rules48}, is the
48-card game (8s removed, eight books, play to the end).

The theorem needs only a handful of rows from the \code{us54} decision table
\cite{rules54}; we reproduce them by number, since the proof cites them.

\begin{center}
\small
\begin{tabular}{@{}c p{0.86\textwidth}@{}}
\toprule
row & rule \\
\midrule
5  & \textbf{Ask target.} The target of an ask is always an opponent; a
     teammate may never be asked. \\
6  & \textbf{Ask licence.} An asker must hold at least one card of the asked
     book. \\
7  & \textbf{Own card.} You may not ask for a card you hold. \\
9  & \textbf{Hit.} The card transfers; the asker keeps the turn and asks again. \\
10 & \textbf{Miss.} The turn passes to the player who was asked. \\
12 & \textbf{Declare content.} Name the book and the exact holder of all six
     cards; every assignment must be a seat on the declarer's own team. \\
14 & \textbf{Any error at all.} The opposing team scores the book. Under
     \code{us54} this covers both ``an opponent held one of the six'' and ``my
     team held all six but I misassigned''; the \code{pagat48} \emph{void}
     outcome is abolished. \\
15 & \textbf{Declare without holding.} Legal --- you may declare a book you
     hold no card of. \\
17 & \textbf{Information.} Card counts are public, with a full public log of
     asks, results, and declares. \\
\bottomrule
\end{tabular}
\end{center}

Under \code{pagat48} the corresponding declare rows split the error case: its
row 14 awards the book to the opposing team when an opponent of the declarer
holds at least one card, regardless of other locations, and its row 15
\emph{voids} the book when the declarer's own team held all six but a location
was wrong \cite{rules48}. The distinction matters in
Section~\ref{sec:correction}. A further \code{us54}-specific rule bears on the
economics: declares may be made at any time, including out of turn, so a
teammate can beat you to a book you were saving \cite{rules54}.

\section{The absorbing-containment theorem}
\label{sec:theorem}

\begin{definition}
A book $B$ is \emph{contained} by team $T$ in a state $s$ if $B$ is unresolved
and every one of its six cards is held by a seat of $T$.
\end{definition}

\begin{theorem}[Absorbing containment]
\label{thm:absorb}
Let $B$ be contained by team $T$. Then, under both \code{us54} and
\code{pagat48}:
\begin{enumerate}
\item[(i)] no seat of the opposing team has a legal ask into $B$;
\item[(ii)] any declare of $B$ by a seat of the opposing team awards $B$ to
$T$; and
\item[(iii)] containment persists until $T$ itself resolves $B$.
\end{enumerate}
\end{theorem}

\begin{proofenv}
(i) Asking into $B$ requires holding at least one card of $B$ (row 6, identical
in both rule sets); an opponent holds none.

(ii) A declare must assign all six cards to seats on the declarer's own team
(row 12). Since every card of $B$ sits with $T$, every assignment made by an
opposing declarer is wrong. Under \code{us54}, row 14 awards a declare with any
error to the opposing team --- which is $T$. Under \code{pagat48} the rule that
fires is likewise its row 14: at least one card of $B$ is held by an opponent
of the declarer (indeed all six are), so the opposing team --- again $T$ ---
scores the book. The \code{pagat48} void (its row 15) cannot apply: it requires
the \emph{declarer's own} team to hold all six, which contradicts containment.

(iii) The cards of $B$ can move only through a hit on an ask into $B$ or
through resolution of $B$. By (i) opponents cannot ask into $B$ at all. A
holder's ask into $B$ must name a card the holder lacks (row 7) at an opponent
--- who holds no card of $B$ --- so it can only miss, and a miss moves no
cards (row 10). Hence every card of $B$ stays with $T$ until $B$ is resolved,
and by (ii) an opponent's attempt to resolve it awards it to $T$. The only
transition out of the contained state that does not score for $T$ is
impossible; the state is absorbing.
\end{proofenv}

\begin{corollary}[The turn-pass]
\label{cor:pass}
Let a seat of $T$ hold $k$ cards of a contained $B$, $1 \le k \le 5$. Then that
seat may legally ask any chosen opponent for any card of $B$ it does not hold
(rows 5--7); the ask is a guaranteed miss, so the turn passes to the chosen
opponent (row 10); the miss moves no cards, so the licence is never consumed
and the move repeats indefinitely; and resolving $B$ destroys the move.
\end{corollary}

These are claims C1--C6 of the project's containment report \cite{containment},
and their evidence tier there is \emph{measured}, not \emph{derived}: each is
executed against the engine in \code{tests/engine/containment.test.ts} on every
commit. In the pinned position (team A at seats $0,2,4$ containing LOW-C), the
opponents have zero legal asks into the book (C1); an opponent's declare
resolves with outcome \code{team0}, the book gifted back (C2); the holder's ask
is legal, misses, and sends turn $0 \to 3$ as aimed (C3/C4); five consecutive
uses leave the hand unchanged (C5); and after claiming, zero asks into the book
remain (C6) \cite{containment}. The independent attestation is Develin's
strategy text \cite{develin}, which advises holding a fully controlled suit
without declaring --- the only way the opponents can win it is if you declare
it wrong --- and separately names asking into it as a way to pass information
to teammates. The project's tier discipline explicitly refuses to count the
recent cluster of AI-generated Literature repositories, including its own
siblings, as corroboration \cite{containment}.

\subsection{The correction: \code{pagat48} is absorbing too, and a draft said otherwise}
\label{sec:correction}

An earlier draft of the project's own documentation claimed that containment
was \emph{not} absorbing under \code{pagat48}, on the argument that its row 15
would void an opponent's declare of a contained book \cite{styles}. That is
false, and it was caught by measurement rather than by re-reading: row 15's
void requires the declarer's own team to hold all six cards, which cannot
describe a seat whose opponents hold all six; the rule that actually fires is
row 14, worded identically in both rule sets. Executed under \code{pagat48},
the opponent's declare of a contained book resolves with outcome \code{team0}
and score $[1,0]$ --- gifted back, not voided --- and the test suite now pins
both C1 and C2 under the 48-card rules so the corrected fact cannot decay into
folklore \cite{styles}.\footnote{Prose decays faster than assertions. At the
time of writing, one bot-level test in
\code{tests/bots/contained.test.ts} still carries the superseded rationale in
its name and comment (``row 14 gifts, row 15 voids'') while asserting behaviour
that remains correct for a different reason; the engine-level suite carries the
correction. The episode is a working example of why this project records
evidence tiers and executes its claims.}

The correction has a practical edge: the policy of
Section~\ref{sec:discipline} refuses to run under \code{pagat48}, and the
refusal is now honestly labelled a \emph{compatibility} gate --- the project
holds the shipped 48-card game byte-identical, verified by differential digest
--- rather than a rules gate. The mechanism is valid there and is declined
anyway; no rules argument stands in the way of revisiting it if the 48-card
roster is ever re-tuned \cite{styles}.

\section{The move and its information price}
\label{sec:move}

The incoming research handoff described the turn-pass as information-free.
Measured, it is not --- but the qualification is narrow and makes the move
better, not worse \cite{containment}. Each ask publishes, via the public log
(row 17), that the asker holds at least one card of the book and \emph{lacks
the named card}. The first time a given card is named, the asker's hand is
narrowed by one card. Naming the \emph{same} card again publishes nothing new,
and the miss is equally guaranteed. The discipline that follows is the
report's own: \emph{reuse one card as your turn-pass}. The first use spends
one card of hand information; every repeat is genuinely free; cycling through
distinct cards leaks a bit each time for no gain. In the measured position the
holder had four distinct nameable cards --- four lifetimes of the move, if it
were ever worth spending them \cite{containment}.

Two framings of the same ask coexist and are deliberately kept apart
\cite{containment,styles}. The rules-derivation frames it as \emph{turn
control}: the miss chooses who acts next. Develin frames it as
\emph{signalling}: the publication reaches the teammates, who --- unlike the
opponents, whom C1 and C2 lock out --- can convert it, for instance by
declaring the book, which C2 denies the other side. Both are modelled; only
turn control decides (Section~\ref{sec:trigger}); signalling enters at exactly
one term, the information price, which a \code{signalling} style books as
delivered rather than spent.

Under \code{us54} the theorem's ``no defensive value in claiming'' meets one
genuine countervailing cost that \code{pagat48} does not share: declares are
legal at any time, so \emph{a teammate may declare the contained book wrongly
first}. That intra-team race --- measured by the roster's \code{raceLosses}
diagnostic --- is the real price of waiting, and it binds the declare policy,
not the ask policy \cite{containment}.

\section{The trigger, derived}
\label{sec:trigger}

Passing the turn is normally bad: you surrender tempo for nothing. The move is
worth playing only when the ordinary alternative is poor \emph{and} aiming the
concession buys something. Both quantities are read off public data (row 17),
in units of cards \cite{styles}.

Write $p^{*}$ for the hit probability of the ask the style would otherwise
play, $t_{o}$ for the opponent that ask concedes the turn to on a miss, and
$t_{c}$ for the opponent the style's aim (\code{missTarget}) would pick. Let
$V_{\mathrm{hit}} = 1$ card (the card taken; the asker also keeps the turn,
row 9) and let $V_{\mathrm{miss}}(t) = -\,\mathrm{cost}(t)$, the cards a
conceded turn is expected to hand seat $t$. Then
\begin{align*}
\mathrm{value}(\text{ordinary}) &= p^{*}\,V_{\mathrm{hit}}
   + (1-p^{*})\,V_{\mathrm{miss}}(t_{o}),\\
\mathrm{value}(\text{pass}) &= V_{\mathrm{miss}}(t_{c}) - \mathrm{infoCost},
\end{align*}
and the pass wins exactly when
\begin{equation}
\label{eq:trigger}
p^{*} \;<\; \frac{a \cdot \mathrm{gain} - \mathrm{infoCost}}{\mathrm{tempo}},
\end{equation}
with every term defined and sourced as follows \cite{styles}:

\begin{center}
\small
\begin{tabular}{@{}l l p{0.55\textwidth}@{}}
\toprule
term & value & where it comes from \\
\midrule
$E$ & $\mathrm{hits}/\max(1,\mathrm{misses})$ &
  over the public log. Row 9 keeps the turn on a hit and row 10 ends it on a
  miss, so a turn is a geometric run of hits with expected length $h/(1-h)$
  --- on counts, exactly hits over misses. \\
$\mathrm{cost}(t)$ & $E \cdot n_{t}/\bar{n}$ &
  a conceded turn is worth more to a seat holding more ask licences (row 6);
  hand size $n_{t}$ is the public part of that (row 17). \\
$\mathrm{gain}$ & $E \cdot (n_{o}-n_{c})/\bar{n}$ &
  what aiming the concession buys: $\mathrm{cost}(t_o)-\mathrm{cost}(t_c)$. \\
$\mathrm{tempo}$ & $1 + E \cdot n_{o}/\bar{n}$ &
  the swing between hitting and missing: one card taken, plus the turn
  \emph{not} conceded to $t_o$. \\
$\mathrm{infoCost}$ & $0$ on a reuse, else $1/U$ &
  the Section~\ref{sec:move} price, charged at the only channel C1 and C2
  leave the opponents: count exhaustion. \\
$a$ & the style's \code{containedPass} &
  the appetite (Section~\ref{sec:appetite}). \\
\bottomrule
\end{tabular}
\end{center}

\noindent
Here $\bar{n}$ is the mean hand size over seats still holding cards and $U$
the number of live cards whose holder is not yet publicly certain. Every input
is public and none is a tuned constant: an engine, a roster, or a rule set
that hits more often prices a conceded turn higher, automatically. $p^{*}$ is
the seat's own \emph{best} estimate --- the constraint-refined probability
where the skill model provides one, not the slot prior --- so the ask being
displaced is never undervalued. Two consequences fall out of the algebra, and
they are why this is not ``fire whenever legal'' \cite{styles}:

\begin{itemize}
\item $\mathrm{gain} \le 0$ whenever the ordinary ask already concedes where
the style wants it, and the move is then refused however contained the book
is. This switches the mechanism off wholesale for two styles with no bespoke
code: Blitz aims at the \emph{largest} hand (\code{missTarget} \code{'most'}),
so its ordinary concession is already its preferred one; a style with
\code{missTarget} \code{'random'} expresses no aim and buys nothing by aiming.
\item A certain hit is never displaced: it is riskless material \emph{and}
keeps the turn (row 9). The threshold is clamped strictly below 1 and the
check is explicit.
\end{itemize}

The implementation is deliberately small enough to quote. The valuation in
\code{lib/engine/bots/contained.ts} is, in its entirety of consequence (lightly
adapted: the file computes these same expressions and returns \code{threshold}
as a field of its result object rather than assigning it):

\begin{quote}
\begin{verbatim}
const E = hits / Math.max(1, misses)
const meanHand = seats > 0 ? sum / seats : 0
const perCard = meanHand > 0 ? E / meanHand : 0
const gain  = perCard * (nOrd - nChosen)
const tempo = 1 + perCard * nOrd
const raw = (style.containedPass * gain - infoCost) / tempo
threshold = Math.max(0, Math.min(1 - 1e-9, raw))
\end{verbatim}
\end{quote}

One knob is deliberately \emph{not} consulted: \code{minHitP}, the roster's
long-shot filter. That knob refuses asks whose chance of taking a card does
not justify the turn they risk; this ask has no chance of taking a card by
construction and risks nothing, because the turn is the entire content of the
move rather than its downside. Applying a hit-probability floor to it would
refuse it always, at every style that carries one \cite{contained-ts}.

A worked instance, pinned by test \cite{tests-bots}: seat 0 holds five of
LOW-C with the sixth placed on a teammate by three public misses, opponents at
$20/5/2$ cards, and the log showing 3 hits to 3 misses, so $E = 1$ and
$\bar{n} = 49/6$. Then $\mathrm{gain} = 18/\bar n \approx 2.20$,
$\mathrm{tempo} = 1 + 20/\bar n \approx 3.45$, threshold $\approx 0.64$: the
pass fires. Add six filler misses so $E = 1/3$ and the same position prices at
$0.404$ for appetite 1 --- below the ordinary ask's $p^{*} = 0.465$, so
Balanced keeps its material ask --- while the Hoarder's appetite of $1.33$
lifts the bar to $0.538$ and it spends the turn to aim the concession at the
two-card seat. The same knob, the same arithmetic, two different styles.

\subsection{The appetite: one number, eight of nine styles share it}
\label{sec:appetite}

Inequality~\eqref{eq:trigger} prices a \emph{single} use, but the licence is
renewable (C5): a style that will hold the book longer gets more uses out of
the same first-use cost. The style's sole contribution is therefore
\code{containedPass}, the \emph{expected number of uses the licence gets
before the book is banked}. A style that banks at its next opportunity takes
$1$; eight of the nine roster styles do. The Hoarder takes $1.33$, read off
the project's own measurement of how much longer it holds a set ---
\code{declareLatency} $31.02$ window-cycles against Balanced's $23.39$ --- and
it is the only appetite above break-even because it is the only delay that has
actually been measured \cite{styles}.

The uniformity is deliberate: which styles \emph{reach} contained books, and
whether aiming buys anything, are already governed by knobs the roster carried
before (\code{hoardBooks}, \code{declareOnlyOwnHand}, \code{missTarget}), and
nine per-style appetites would have hidden that behind nine tuned constants.
An honest check on the number: measured uses per contained book run Turtle
$3.40$, Hoarder $2.64$, Banker $2.56$, Balanced $2.23$, Punter $1.63$ --- the
\emph{ordering} the appetite argument predicts, but a Hoarder-to-Balanced
ratio of $1.18$, not the $1.33$ the latency ratio implied. The number is right
in sign and roughly right in size, and it has deliberately not been re-derived
from $1.18$, because fitting an appetite to the behaviour it produces is
tuning against the outcome \cite{styles}.

\subsection{Implementation discipline}
\label{sec:discipline}

Three lines the implementation refuses to cross, each pinned by test rather
than asserted \cite{contained-ts,tests-bots}:

\textbf{Card reuse, stateless.} The card chooser prefers a card of the book
whose absence from this hand is already public; failing that it takes the
canonical-first card the seat does not hold --- which, because a contained
book's cards can never move (Theorem~\ref{thm:absorb}(iii)), is fixed for the
rest of the game, so every later call finds that same card already in the log.
The reuse discipline holds with no state carried between decisions. Measured
in play, $97.6\%$ of the Turtle's uses and $94.8\%$ of the Hoarder's cost no
information at all \cite{styles}.

\textbf{No partner model.} Develin records that prearranged conventions are
forbidden in this tradition, and the containment report says not to blur the
line \cite{containment,develin}. Signalling through a legal public ask is not
a prearranged convention --- but only as long as nothing decodes it. So
nothing does: a teammate's knowledge state after a turn-pass is bit-for-bit
the knowledge it builds after any other miss on the same card by the same seat
--- the row-17 public facts and nothing else --- and the test suite asserts
that equality directly, comparing the two knowledge objects. The inference
layer does not know the ask was a turn-pass, and must not be taught to.

\textbf{Gates before arithmetic.} The plan is refused, cheapest check first,
on: a zero appetite (every shipped difficulty tier has \code{containedPass}
$= 0$, keeping the live product untouched); the rule set (the
Section~\ref{sec:correction} compatibility gate); a skill that cannot read
hand counts (a seat that cannot compare hand sizes cannot aim); a certain hit;
then the recogniser; then Inequality~\eqref{eq:trigger}. The recogniser's
predicate is deliberately weaker than the report's measured position: it
requires every card of the book to be \emph{certainly on the seat's team},
not pinned to named seats --- a pinned book is one the certain-claim path has
already banked, so a pinning recogniser would fire only for the two styles
that refuse certain declares. Under the weaker predicate the ask still cannot
hit, and every style can reach the state \cite{contained-ts}.

\section{Measurements}
\label{sec:measure}

\subsection{Instruments}
\label{sec:instruments}

Four instruments, in increasing order of scale, following the lab's
duplicate-deal and fixed-precision protocol \cite{styles,botlab}:
(i)~\emph{decision mirrors}: 100 \code{us54} games per style in which two
vectors identical but for \code{containedPass} are handed the same seat view
and seed at every decision point and their chosen actions compared;
(ii)~\emph{duplicate-deal pairs}: 1{,}200 pairs (2{,}400 games) per style,
both orientations of every seed with a shared seat rotation, the pair as the
unit of analysis; (iii)~\emph{matrix v2}: the full $36$-cell round-robin at
4{,}300 pairs per cell (309{,}600 games, standard error $\le 0.005$ per cell)
re-run on the same seed set as matrix v1, so the two runs differ by the policy
change and nothing else --- both artifacts committed
(\code{src/lab/data/style-results.dominant.json}, records digest
\code{087b626cedc5845e}, engine commit recorded \code{819eebb-dirty};
\code{style-results.v2.json}, digest \code{0c3cb524a3534d4a}, engine commit
recorded \code{1667a1d-dirty} --- the \code{-dirty} suffixes reflecting
generation from a modified working tree at those commits); and
(iv)~\emph{differential digests}: seeded games
serialised action-by-action and hashed against the prior commit. The v2 health
gate is clean: 0 illegal actions, 0 capped games, 0 invariant violations, 0
ties, 0 voids, 0 non-clinch terminations, 4{,}300 of 4{,}300 distinct seeds.

The digests bound the blast radius exactly \cite{styles}: under
\code{pagat48}, both the three shipped tiers and all nine roster styles are
bit-identical to the prior commit; under \code{us54} the shipped tiers are
bit-identical (every preset carries \code{containedPass} $=0$) and the roster
changed for seven of nine styles --- the two unchanged, Blitz and Scout, being
exactly the two the model predicts (aim already at the largest hand; almost
never holds a contained book). Even that count needed correcting by
measurement: at 40 games the Archivist also digested as unchanged, because 40
games is below its firing rate; at 250 duplicate pairs (500 games) its digest
differs, and 7-of-9 is the corrected figure.

\subsection{Opportunity, fire, and refusal}
\label{sec:fire}

Table~\ref{tab:fire} gives the decision-mirror rates. \code{opp} counts ask
decisions at which a contained book was available at all; \code{fire} counts
decisions at which the valuation returned a plan; in every style the number of
divergent decisions equals \code{fire} exactly --- the mechanism never once
chose the move the ordinary policy was already going to play \cite{styles}.

\begin{table}[ht]
\centering
\caption{Opportunity and fire rates, 100 mirror \code{us54} games per style
\cite{styles}.}
\label{tab:fire}
\small
\begin{tabular}{@{}l r l r r r r r r@{}}
\toprule
style & appetite & aim & decisions & opp & opp\,\% & fire & fire/opp & div\,\% \\
\midrule
Balanced  & 1    & fewest & 9{,}741  & 379    & 3.89\%  & 29  & 7.7\%  & 0.0414\% \\
Blitz     & 1    & most   & 9{,}345  & 472    & 5.05\%  & 0   & 0.0\%  & 0.0000\% \\
Punter    & 1    & fewest & 9{,}507  & 347    & 3.65\%  & 13  & 3.8\%  & 0.0190\% \\
Banker    & 1    & fewest & 9{,}818  & 495    & 5.04\%  & 69  & 13.9\% & 0.0981\% \\
Turtle    & 1    & fewest & 17{,}285 & 9{,}310 & 53.86\% & 679 & 7.3\%  & 0.5575\% \\
Hoarder   & 1.33 & fewest & 11{,}418 & 1{,}707 & 14.95\% & 211 & 12.4\% & 0.2580\% \\
Scout     & 1    & fewest & 11{,}292 & 43     & 0.38\%  & 0   & 0.0\%  & 0.0000\% \\
Ghost     & 1    & fewest & 9{,}996  & 131    & 1.31\%  & 21  & 16.0\% & 0.0293\% \\
Archivist & 1    & fewest & 10{,}383 & 105    & 1.01\%  & 5   & 4.8\%  & 0.0067\% \\
\bottomrule
\end{tabular}
\end{table}

Three readings \cite{styles}. First, \emph{the opportunity is a property of
the declare policy, not the ask policy}: a contained book exists at $53.86\%$
of the Turtle's ask decisions (it never banks a set split across teammates,
so it accumulates them) and $0.38\%$ of the Scout's (whose information
weighting pins holders early, and a pinned contained set is banked
immediately). The Hoarder sits second at $14.95\%$ --- nearly four times the
control's $3.89\%$ --- which is precisely the benefit the project once said it
was paying for and not collecting. Second, \emph{it is mostly a refusal}: even
where the licence exists it is spent at $3.8$--$16.0\%$, and every refusal is
the $\mathrm{gain} \le 0$ arm --- the ordinary ask was already conceding
where the style wanted. Third, \emph{divergence is small}: the largest mover
is the Turtle at $0.56\%$ of all decisions; seven of nine move less than
$0.1\%$; exactly two --- Blitz and Scout --- do not move at all (the
Archivist fires five times, $0.0067\%$). The information price, for completeness:
charged only on first uses ($2.4\%$ of the Turtle's, $5.2\%$ of the
Hoarder's), bounded by $1/U < 0.05$ cards in these positions against a tempo
above 1, and it never flipped a decision.

\subsection{Does it win? No.}
\label{sec:mirror}

Table~\ref{tab:null} is the head-to-head: each style \emph{with} the mechanism
against the identical style with it forced off, 1{,}200 duplicate-deal pairs
each; $0.5$ is no effect; \code{changed} counts games whose result or length
differed from the same deal played by two copies of the unequipped vector ---
games the mechanism actually touched \cite{styles}.

\begin{table}[ht]
\centering
\caption{Per-style null result: WITH vs.\ WITHOUT \code{containedPass},
1{,}200 duplicate-deal pairs (2{,}400 games) per style \cite{styles}.}
\label{tab:null}
\small
\begin{tabular}{@{}l r r r r r@{}}
\toprule
style & appetite & pairs (games) & WITH score & SE & changed \\
\midrule
Balanced  & 1    & 1{,}200 (2{,}400) & 0.5046 & 0.0034 & 239 \\
Blitz     & 1    & 1{,}200 (2{,}400) & 0.5000 & 0.0000 & 0 \\
Punter    & 1    & 1{,}200 (2{,}400) & 0.4971 & 0.0029 & 142 \\
Banker    & 1    & 1{,}200 (2{,}400) & 0.5021 & 0.0035 & 281 \\
Turtle    & 1    & 1{,}200 (2{,}400) & 0.5008 & 0.0065 & 1{,}268 \\
Hoarder   & 1.33 & 1{,}200 (2{,}400) & 0.4938 & 0.0054 & 915 \\
Scout     & 1    & 1{,}200 (2{,}400) & 0.5000 & 0.0000 & 0 \\
Ghost     & 1    & 1{,}200 (2{,}400) & 0.4958 & 0.0027 & 156 \\
Archivist & 1    & 1{,}200 (2{,}400) & 0.4996 & 0.0004 & 11 \\
\bottomrule
\end{tabular}
\end{table}

No style is two standard errors from $0.5$. The largest deviation in either
direction is $1.6$\,SE (Ghost, $0.4958 \pm 0.0027$); the two heaviest users
land on and slightly below the control (Turtle $0.5008 \pm 0.0065$, Hoarder
$0.4938 \pm 0.0054$); across the seven styles that fire at all, point
estimates straddle $0.5$ (four below, three above, mean $0.4991$). Blitz and
Scout return exactly $0.5000 \pm 0.0000$ over 2{,}400 games each, which is the
zero-changed-games column restated: for those two, the vectors are the same
player \cite{styles}.

\subsection{The full matrix confirms it}
\label{sec:matrix}

The mirror isolates the mechanism against a twin; matrix v2 exposes it to the
whole roster. Nothing that matters moved \cite{styles}: the verdict is
unchanged (dominant style: Punter, all four criteria); the ranking is
identical position for position; Punter's mean score is $0.5829$ in both runs
and its maximin $0.5193 \to 0.5190$ (its worst matchup relabelled between two
cells that were already tied); cyclic energy rose $0.0097 \to 0.0112$ against
a threshold of $0.15$, with the same single directed 3-cycle
(Hoarder $>$ Archivist $>$ Ghost) still failing Benjamini--Hochberg
($q\ 0.31 \to 0.27$); significant cells went $34 \to 33$ of 36, entirely
because one marginal Hoarder cell fell out of significance; Nash and
$\alpha$-Rank both still put all mass on Punter. No cell moved significantly:
the largest paired per-deal difference is \code{blitz-vs-ghost} at
$+0.0030 \pm 0.0015$ ($t = 2.07$), which does not survive
Benjamini--Hochberg over 36 comparisons --- one $|t|$ above 2 in 36 is what
chance produces. Per style, paired over the same 4{,}300 deals, every mean
score is within $1.5$\,SE of zero change (Table~\ref{tab:v2}); the headline
row is the Hoarder --- the style built to exploit the mechanism --- at
$\Delta = -0.0018 \pm 0.0014$ ($t = -1.35$), sixth of nine in both runs.

\begin{table}[ht]
\centering
\caption{Matrix v2 vs.\ v1: per-style mean score, paired over the same
4{,}300 deals per cell \cite{styles,artifacts}.}
\label{tab:v2}
\small
\begin{tabular}{@{}l r r r r r l@{}}
\toprule
style & v1 & v2 & $\Delta$ & SE($\Delta$) & $t$ & rank \\
\midrule
Punter    & 0.5829 & 0.5829 & $-0.0001$ & 0.0009 & $-0.07$ & 1 $\to$ 1 \\
Blitz     & 0.5594 & 0.5602 & $+0.0008$ & 0.0008 & $+0.94$ & 2 $\to$ 2 \\
Balanced  & 0.5584 & 0.5581 & $-0.0003$ & 0.0009 & $-0.34$ & 3 $\to$ 3 \\
Banker    & 0.5274 & 0.5285 & $+0.0011$ & 0.0011 & $+1.00$ & 4 $\to$ 4 \\
Ghost     & 0.4928 & 0.4928 & $+0.0000$ & 0.0010 & $+0.03$ & 5 $\to$ 5 \\
Hoarder   & 0.4904 & 0.4885 & $-0.0018$ & 0.0014 & $-1.35$ & 6 $\to$ 6 \\
Archivist & 0.4791 & 0.4802 & $+0.0011$ & 0.0008 & $+1.47$ & 7 $\to$ 7 \\
Scout     & 0.4064 & 0.4063 & $-0.0001$ & 0.0007 & $-0.12$ & 8 $\to$ 8 \\
Turtle    & 0.4032 & 0.4025 & $-0.0007$ & 0.0019 & $-0.38$ & 9 $\to$ 9 \\
\bottomrule
\end{tabular}
\end{table}

Two rows of Table~\ref{tab:v2} are not evidence of anything and should not be
read as such: Blitz's and Scout's deltas are nonzero even though their own
policies are byte-identical between runs, because their \emph{opponents}
changed \cite{styles}.

\paragraph{Which v2 this is (note added August 2026).} The v2 column above is
matrix v2 as generated at engine commit \code{1667a1d}, records digest
\code{0c3cb524a3534d4a} \cite{artifacts}. \textbf{The comparison it makes is
unaffected by anything that has happened since}, and for the reason that makes
it a comparison at all: v1 and that v2 were run on the same seed set and
differ by the turn-pass and nothing else, so the paired deltas isolate the
mechanism. The artifact \emph{path} has since been rewritten by a controlled
re-measurement of the same 36 cells --- same seeds, same sample size --- run
against a roster that additionally carries the concession layer, and with a
larger exploitability ladder. The verdict survives it (Punter still dominant
on all four criteria) but every headline figure moved. That artifact differs
from v1 by two changes rather than one and therefore \emph{cannot} be
substituted into Table~\ref{tab:v2}: doing so would measure the turn-pass and
the concession layer together and attribute the sum to the turn-pass. Where a
current reading of the roster is what is wanted --- ranks, mean scores,
exploitability --- the superseding numbers are in the v0.5 paper's addendum
\cite{v05}, and they are not restated here.

\subsection{Exploitability, and the axis the search cannot turn}
\label{sec:exploit}

The single-knob exploitability search was re-run from scratch for all nine
styles (its cache keys on a hash of the engine tree, so the new file
invalidated every entry). The deltas are all inside the search's own noise
--- its final measurements carry standard errors of $0.008$--$0.017$ and it
only accepts moves worth at least $\approx 0.03$ --- and the verdict criterion
passes in both runs, with the dominant style's exploitability moving
$0.0000 \to 0.0025$ against a rivals' median of $0.0444$ \cite{styles}. The
one qualitatively interesting row is the Hoarder's: the best single-knob
counter to the Hoarder-with-the-benefit is no longer a scoring weight but
\code{hoardBooks=0 minHandSize=0} --- \emph{stop hoarding} --- worth
$0.5563 \pm 0.0124$. The counter to the style is still to not be it
\cite{styles,artifacts}.

One gap remained: the search's knob ladder predates \code{containedPass}, so
neither run could switch the mechanism itself on or off. That axis was probed
directly with the same sequential protocol used for tuning (SPRT,
$d_0 = 0$, $d_1 = 0.03$, $\alpha = \beta = 0.05$, 24--400 pairs, then a
fixed-$N$ 400-pair evaluation on unseen seeds), against every style, at
appetites $\{0, 2, 4\}$. \textbf{No candidate was accepted, for any style, at
any appetite.} The two rows worth naming: turning the mechanism \emph{off}
against the Turtle ran to the pair cap unresolved
($+0.0125 \pm 0.0120$ in search; $0.5225 \pm 0.0121$ in eval) --- the one
candidate the SPRT could not decide; and the same candidate against the
Hoarder produced a sign disagreement between search and eval
($-0.0200 \pm 0.0176$ rejected, versus $0.5275 \pm 0.0098$), recorded
precisely because it is the one number in the exercise that would have looked
like a finding if only its second half had been reported --- the honest
summary of two contradictory $0.01$--$0.02$-sized effects is that neither is
real. What \emph{does} resolve is over-use: at appetite 4 the Turtle drops to
$0.4600 \pm 0.0148$ and the Banker to $0.4850 \pm 0.0083$. The move is worth
nothing, and spending turns on it anyway is worth less than nothing
\cite{styles}.

\paragraph{The ladder gap has since been closed (note added August 2026).}
\code{containedPass} is now a rung of the exploitability search itself, along
with two knobs of the concession layer, and the re-measurement noted in
Section~\ref{sec:matrix} re-ran every style against the enlarged ladder. Every
$E$ in this subsection --- the dominant style's $0.0000 \to 0.0025$, the
rivals' median of $0.0444$, and the Hoarder's \code{hoardBooks=0
minHandSize=0} --- is a figure of the old ladder and is superseded. The new
values, and the reason none of their magnitudes may be attributed to the
ladder alone, are reported in the v0.5 paper's addendum \cite{v05}; they are
not restated here. \emph{This paper's conclusion does not turn on them}: the
mechanism's null is measured by the mirror of Table~\ref{tab:null} and by the
SPRT probe above, and neither reads the payoff artifact. The direct probe of
the \code{containedPass} axis reported here also stands as its own record ---
it is what the search could not reach at the time, measured by other means,
and its answer and the enlarged ladder's agree that the axis is not where an
exploit lives.

\section{Why the effect is small}
\label{sec:small}

The null is not mysterious; the trigger's own terms bound it. A conceded turn
is worth $E \cdot n_t/\bar{n}$ cards, and the measured $E$ --- hits over
misses on the public log --- is around 1 at this roster's level of play. So
the \emph{entire} spread between the best and the worst opponent to concede
to is on the order of one card, and the aimable fraction of that spread is
all the move can ever buy. Against it stands everything the ordinary ask
gives up: $p^{*}$ of a card of material, plus the turn a hit would have
retained (row 9). The move is real, and it is small \cite{styles}. The
firing profile of Table~\ref{tab:fire} says the same thing from the other
side: most of the time the concession is already aimed correctly, so there is
nothing to buy, and the valuation --- honestly computed --- mostly refuses.

Under \code{us54} there is also an active cost to the holding pattern the
move requires: the intra-team race. A contained book left unclaimed is a book
a teammate can declare wrongly at any time, gifting it away (rows 11 and 14
of \cite{rules54}); the Hoarder's \code{raceLosses} doubles ($0.046 \to
0.091$ per game) when its hoarding is actually expressed \cite{styles}. The
absorbing state is safe from the opponents, not from one's own side.

\section{What is not settled}
\label{sec:open}

\textbf{The declare-policy mechanism.} Everything measured here changes the
\emph{ask} policy: when a style already holds a contained book, it may spend
a turn aiming a concession. The containment report's second recommendation
--- hold contained books by default, trading tempo and race risk against the
retained licence, instead of treating banking as automatically correct ---
is a change to the \emph{declare} policy and is untouched
\cite{containment,styles}. Table~\ref{tab:fire}'s first reading is the
evidence that this is where the leverage would be: the opportunity rate is
set almost entirely by how long a style leaves books unclaimed ($53.86\%$
for the Turtle against $0.38\%$ for the Scout), i.e.\ by the declare policy.
The report's expected-value analysis is on file for whoever attempts it:
with $c = P(\text{contained})$ and
$a = P(\text{assignment correct} \mid \text{contained})$, declare iff
$c > 1/(1+a)$ under \code{pagat48}'s void rule and $c \cdot a > 1/2$ under
\code{us54}'s gift rule --- and do not collapse $c$ and $a$ into one number,
as the current claim planner does, since ``is this book on my team'' and
``do I know which teammate holds each card'' have different distributions
and different remedies, and conflating them biases the bot conservative
\cite{containment}.

\textbf{An untuned roster.} The bar Inequality~\eqref{eq:trigger} must clear
is \emph{the best ordinary ask the style has}, and the project's tuning
protocol has never been run over this roster; a roster with a weaker ask
policy would leave more room for a turn-pass \cite{styles}. Relatedly, the
deliberate known-miss ask has been independently flagged as a gap: the
FishBot v0.3 paper's engine also permits known-miss asks and its policy
``almost never'' chooses them \cite{containment}. Two independent sources
plus a measured mechanism made the gap real; this paper measures what closing
it is worth at one operating point, not at all of them.

\section{Threats to validity}
\label{sec:threats}

\emph{The claims are about this rule set as implemented.} C1--C6 are
measured against this engine; that is the strongest tier available here, and
also its scope. The independent attestation (Develin, read directly) covers
the absorbing property and both framings of the ask, not the quantitative
results \cite{containment}.

\emph{Self-play, one roster, one engine.} The null is measured between
policy vectors sharing one inference engine, on a nine-style roster at one
level of play. A stronger or weaker population could price the move
differently; Section~\ref{sec:open} names the direction.

\emph{The appetite is a proxy.} \code{containedPass} $= 1.33$ is read off a
latency ratio, and the realised use ratio is $1.18$. The discrepancy is
reported and deliberately left unfitted \cite{styles}.

\emph{Search noise.} The exploitability deltas of Section~\ref{sec:exploit}
sit inside the search's own error bars, and the one sign-contradictory
result is reported as noise, not evidence. The SPRT probe of the
\code{containedPass} axis closes the ladder gap but inherits the same
detection floor ($\approx 0.03$); effects smaller than that would not have
been accepted --- though the paired mirror of Table~\ref{tab:null}, with
standard errors of $0.003$--$0.007$, would have seen them.

\emph{Small-sample traps, twice.} Forty games misread the Archivist's digest
as unchanged; a 200-pair pilot misplaced the Hoarder's rank (7th, against 6th
in both full-precision runs). Both errors were caught by re-measuring at
scale and both corrections are recorded in the source documents
\cite{styles}. They are kept in this paper as calibration for how much to
trust any single small run --- including the ones reported here.

\section{Conclusion}
\label{sec:conclusion}

A theorem can be true, its mechanism implementable, its implementation
faithful, and the whole thing worth nothing. The contained book is absorbing
under both rule sets this project supports --- a fact strong enough to
survive the project's own attempt to disprove half of it --- and the
turn-pass it licenses is renewable, aimable, and after its first use free.
Implemented with an honest price and handed to the style built to exploit
it, it changed 915 of that style's 2{,}400 mirror games and moved its score
by $-0.0062 \pm 0.0054$ against its twin and $-0.0018 \pm 0.0014$ against
the field. The measured value of the move at this roster's level of play is
zero, and the arithmetic that prices it says why: a conceded turn is worth
about one card here, and one card, aimed or not, does not win races to five.

The negative result is the deliverable. It converts ``the Hoarder
underperforms because its benefit was never implemented'' --- a live
objection, on the record --- into ``the benefit is implemented, reached
nearly four times as often as baseline, priced above break-even, and worth
nothing,'' and it does so with the objection's own instruments. What remains
open is exactly one mechanism (the declare-side holding policy), and the
measurements here say where its leverage would have to come from. Measuring
rather than asserting cost two corrections on the way to a zero; both
corrections, and the zero, are the point.

\section{Addendum: the artifact this paper pairs, and the one the lab now serves}
\label{sec:addregen}

\emph{September 2026.} Two things happened to the matrix after this paper was written, and
neither disturbs its result --- but the second changes what a reader will find if they open the
lab looking for these numbers, so it is recorded here rather than left to be discovered.

\textbf{First, a rules correction.} \code{RULES\_US54.md} \S4 marked a rule \code{[DERIVED]} and
the inference was wrong: when a declare emptied the current turn-holder the engine advanced the
turn to the next seat in ascending cyclic order holding cards, which is an opponent whenever that
opponent holds a card, where the owner's rulebook passes the turn to a teammate. Over 400 games,
19 of the 56 such events handed the turn to the wrong team \cite{rules54}. The engine is
corrected and the whole corpus was regenerated at the new rules hash.

\textbf{Second, and consequently, \code{style-results.v2.json} has been regenerated again.} The
file path this paper's Section~\ref{sec:matrix} names now holds a fourth measurement of that
matrix, not the one Table~\ref{tab:v2} reports.

\subsection{Why the central result is untouched}

\textbf{Nothing in Sections~\ref{sec:mirror}--\ref{sec:small} may be re-derived from the current
artifact, and nothing in them needs to be.} This paper's measurement of the contained-book
turn-pass is a \emph{paired} comparison between two committed artifacts, v1 and v2, which differ
by the turn-pass policy and by nothing else \cite{artifacts}. That is what makes
$\Delta = -0.0018 \pm 0.0014$ a measurement of the move rather than of an engine generation. The
artifact now on that path differs from v1 by the turn-pass, the concession layer \emph{and} the
rules correction, so substituting it would measure three changes and attribute them to one. The
v1 and v2 files this paper reads are unchanged on disk; the pairing stands exactly as printed.

\subsection{What the correction was worth, where it can be seen}

The rules correction is measurable on the matrix as a whole and it is very small. Re-running the
identical 36 cells and 4{,}300 duplicate pairs with the turn-pass rule as the only difference
moves no cell by as much as two of its own standard errors: the mean absolute per-cell move is
$0.0016$ and the largest is $0.0038$, against a per-cell standard error of about $0.005$. Punter's
mean score moves by $0.0006$ and the dominance verdict is unchanged \cite{v05}.

\textbf{That is a fact about the instrument, not a verdict on the defect.} A duplicate-deal design
pairs two arms that both carry the rule, so a perturbation symmetric between them cancels --- and
this one was symmetric. A defect firing in about one game in twenty-one was invisible to a
309{,}600-game measurement. Section~\ref{sec:small}'s argument that a conceded turn is worth about
one card, and that one card does not win races to five, is if anything strengthened by that: this
matrix cannot resolve effects of the size a single turn is worth, whoever the turn goes to.

\begin{thebibliography}{99}

\bibitem{v05}
A.~Hsieh,
\emph{FishAI v0.5: measuring play style without skill in a 54-card Literature
variant},
companion paper, published alongside this one, 2026. The engine, the
nine-style roster, the tournament design, and the full-matrix results this
paper's measurements plug into. Its addendum, ``matrix v2, re-measured'',
reports the August 2026 re-measurement: the verdict survives, every headline
figure moves, and criterion~4 is materially weaker.

\bibitem{containment}
A.~Hsieh,
\emph{CONTAINMENT.md --- the contained-book result},
FishAI project internal report, August 2026.
Claims C1--C6 with evidence tiers, the information-cost correction (\S1.2),
the Develin attestation (\S2), and the claim-EV thresholds (\S3.3).

\bibitem{styles}
A.~Hsieh,
\emph{STYLES.md --- the play-style roster under \code{us54}},
FishAI project internal report, August 2026.
\S6.2 (the Hoarder prediction and its refutation), \S6.3 (the turn-pass
implemented and measured; \S6.3.7 the \code{pagat48} correction), \S6.4
(matrix v2).

\bibitem{rules54}
A.~Hsieh,
\emph{RULES\_US54.md --- the ``US student'' 54-card variant},
FishAI project internal report, August 2026. Decision-table rows cited by
number throughout.

\bibitem{rules48}
FishAI project,
\emph{RULES.md --- the 48-card pagat baseline},
internal report, 2026. Rows 14--15 (the split declare-error outcomes,
including the void).

\bibitem{botlab}
FishAI project,
\emph{BOT\_LAB.md --- simulation-lab protocol},
internal report, 2026. Duplicate-deal pairing (\S5.1--5.2), fixed-$N$ matrix
precision (\S5.4), SPRT tuning (\S5.5), and the exploitability ladder
(\S5.7).

\bibitem{develin}
M.~Develin,
\emph{The Ten Best Card Games You've Never Heard Of}, ch.~9 (Canadian Fish),
\url{http://www.bantha.org/~develin/cardgames.html}.
Independent attestation of the absorbing property and of the signalling
framing; records the tradition's ban on prearranged conventions.

\bibitem{pagat}
J.~McLeod (ed.),
\emph{Literature (Canadian Fish)},
Pagat card-game rules archive,
\url{https://www.pagat.com/quartet/literature.html}.
Folklore source for both rule sets' common core.

\bibitem{contained-ts}
FishAI project,
\code{lib/engine/bots/contained.ts},
the turn-pass policy: recogniser, card reuse, valuation, and gates. The
mechanism claims are pinned by \code{tests/engine/containment.test.ts}
(C1--C6 under \code{us54}; C1--C2 under \code{pagat48}; an end-to-end
policy-through-reducer step).

\bibitem{tests-bots}
FishAI project,
\code{tests/bots/contained.test.ts},
the decision-level pins: the refusal on a level board, the measured-$E$
worked positions, the certain-hit clamp, the appetite's scaling, the
no-partner-model equality, and the shipped tiers' zero appetite.

\bibitem{artifacts}
FishAI project,
committed measurement artifacts
\code{src/lab/data/style-results.dominant.json} (matrix v1, engine commit
recorded \code{819eebb-dirty}, records digest \code{087b626cedc5845e}) and
\code{src/lab/data/style-results.v2.json} (matrix v2, engine commit
recorded \code{1667a1d-dirty}, records digest \code{0c3cb524a3534d4a});
309{,}600 games each, health gates clean. The v1 file is unchanged. The v2
\emph{path} was later overwritten by a controlled re-measurement (engine
commit \code{1fdd22e}, records digest \code{1d7d9f30b59c00f6}, generated
2026-08-31); the digest cited here is the artifact this paper measured
against, recoverable at that path from the repository history, and it is the
one Table~\ref{tab:v2} pairs with v1. See \cite{v05}.

\end{thebibliography}

\end{document}
